\section{Problem 264: at most countably many exceptional factorial perturbations}
\subsection*{1) OBJECTIVE AND SOURCE REVIEW}
Attempt dated 2026-09-23. The supplied \texttt{264.tex}, read completely, investigates whether every well-defined sum $\sum_n1/(a_n+b_n)$ is irrational when $(b_n)$ is a bounded sequence of nonzero integers, for $a_n=2^n$ or $a_n=n!$. Its geometric construction addresses the first sequence negatively; its small-dimension argument does not settle the factorial sequence. I strengthen the latter description by proving eventual uniqueness of the perturbation coding. This is not a proof of the required universal irrationality.
\subsection*{2) WORK}
More generally suppose $a_{n+1}/a_n\to\infty$. Fix $B\ge1$, and restrict each perturbation to the finite alphabet $D_B=\{-B,\ldots,-1,1,\ldots,B\}$. For sufficiently large $m$, two distinct choices at index $m$ give
\[
\left|\frac1{a_m+b}-\frac1{a_m+c}\right|
\ge\frac1{(a_m+B)^2}.
\]
On the other hand the total possible difference of all later terms is at most
\[
\sum_{n>m}\frac{2B}{(a_n-B)^2}.
\]
This is $o(a_m^{-2})$. To see this without any summability assumption beyond the growth condition, take $m$ so large that $a_{n+1}\ge R a_n$ for $n\ge m$ and $a_n\ge2B$. The last display is at most
\[
8B\sum_{j\ge1}a_{m+j}^{-2}
\le\frac{8B}{R^2-1}a_m^{-2}.
\]
Choose a fixed $R$ large enough that this is less than $(a_m+B)^{-2}$ for all sufficiently large $m$.

Thus two tails whose first differing digit occurs sufficiently late cannot have the same sum: the first difference dominates everything after it. Once the initial finite perturbation prefix is fixed, the map from infinite tails to sums is injective. There are only countably many rational numbers, so at most countably many such tails have rational sums. There are finitely many prefixes for each fixed $B$; taking the countable union over $B$ proves that at most countably many bounded nonzero integer perturbation sequences can yield a rational sum. In particular this applies to $a_n=n!$.
\subsection*{3) ADVERSARIAL VERIFICATION}
This countability statement is far weaker than saying there are no rational sums. An injective Cantor-type coding can contain prescribed rational values. Under independent uniform choices from any fixed $D_B$, each individual code has probability zero, so irrationality holds with probability one, but the problem asks about every code. Finite initial denominators equal to zero are excluded; other finite initial terms simply shift the tail by a rational amount and do not affect the argument.
\subsection*{4) FINAL}
\textbf{UNRESOLVED.} Almost every bounded factorial perturbation is irrational, and potential exceptions form a countable family. Excluding even one exceptional sequence is the missing universal step.
\subsection*{5) SOURCES CHECKED}
The complete supplied version; V. Kova\v{c} and T. Tao, \emph{On several irrationality problems for Ahmes series}, \texttt{https://arxiv.org/abs/2406.17593}, current abstract checked 2026-09-23 for context. The coding proof above is self-contained and is not attributed to a theorem not inspected there.
\subsection*{6) COMPLETION ESTIMATE}
Subjective progress toward excluding every factorial exception.
\textbf{COMPLETION: 15\%}
